% Chapter 6 -- Rubric: PO11, PO9 (9.1--9.4), PO6 (6.1), PO7 (7.1), PO8 (8.1--8.4)
\chapter{Project Management, Teamwork, and Impact}
\label{ch:management}

\section{Project Management and Finance}

\subsection{Planning and Risk Management}
\label{sec:plan}
\guidance{Present the plan: phases, milestones, task distribution, and schedule,
noting where actual progress differed. List the major risks and how they were
mitigated.}
\begin{figure}[htbp]
  \centering
  \begin{ganttchart}[
      hgrid, vgrid, x unit=0.9cm,
      time slot format=isodate-yearmonth, time slot unit=month,
      title label font=\small, bar label font=\small,
      bar/.style={fill=gray!40}
    ]{2026-01}{2026-12}
    \gantttitlecalendar{year, month=shortname} \\
    \ganttbar{Requirements}{2026-01}{2026-02} \\
    \ganttbar{Design}{2026-02}{2026-04} \\
    \ganttbar{Implementation}{2026-04}{2026-09} \\
    \ganttbar{Testing}{2026-08}{2026-10} \\
    \ganttbar{Report}{2026-10}{2026-12}
  \end{ganttchart}
  \caption{Project schedule}
  \label{fig:gantt}
\end{figure}

\subsection{Budgeting and Resource Identification}
\label{sec:budget}
\guidance{Identify the resources required (people, hardware, software, cloud
and API services, data) with an itemised budget.}

\subsection{Economic Analysis and Financial Prospecting}
\label{sec:economics}
\guidance{Evaluate the financial prospect of the project with clear
interpretation: development and operating cost, the potential users or market
established in \Cref{sec:market}, and how the solution could generate value or
revenue.}

\subsection{Sustainability in Business and Commercialisation}
\label{sec:commercialisation}
\guidance{Explain how the solution could remain viable after the course:
the commercialisation or deployment path (e.g. with the industry partner) and
who would maintain and operate it.}

\section{Individual and Teamwork}
\label{sec:team}

\subsection{Participation and Contribution}
\label{sec:participation}
\guidance{For each member, give the role, specific contributions, participation
in discussions and decision-making, and whether responsibilities were completed
on time.}

\subsection{Collaboration and Conflict Resolution}
\label{sec:collaboration}
\guidance{Describe how the team collaborated and kept a constructive atmosphere,
and how any conflicts were identified and resolved.}

\subsection{Leadership and Direction}
\label{sec:leadership}
\guidance{Explain how direction toward the goals was provided, how members' views
were listened to, how the team planned for improvement after setbacks, and how
members were motivated.}

\subsection{Multidisciplinary Engagement}
\label{sec:multidisciplinary}
\guidance{Explain the multidisciplinary nature of the project and the team's
participation in activities outside its own discipline: work with the industry
partner, domain experts, end users, or other departments.}

\section{Impact on Society}
\label{sec:society}

This system does not reach the public directly. It tests other people's
software, so its effect on society travels through two channels: the quality of
the applications that reach users, and the work of the people who currently
test them. Both are considered below, followed by the negative impacts and what
the design does about them.

\paragraph{Software quality, and why it is a safety question here.} A defect in
the graphical interface of a consumer application is usually an annoyance. In
the class of application this project was evaluated against it is not. The
target is an agricultural commerce platform used by smallholder farmers and
buyers in Bangladesh, where a wrong price, an order routed to the wrong seller,
or a total that does not match its line items is a direct financial loss to a
party with very little margin to absorb it. Those are exactly the defects that
a crash detector cannot see, for the reason set out in
\Cref{subsec:measurement}, and exactly the class that a specification-grounded
oracle is built to catch. \emph{Who is affected:} the end users of any
application tested with this system. \emph{How the design responds:} the oracle
is derived from the requirements rather than from the absence of a crash, and
every generated test cites the requirement it covers, so a defect report can be
traced to the behaviour that was promised.

\paragraph{Access to systematic testing.} Cost is not a neutral engineering
detail in this context. A firm that cannot fund a dedicated quality-assurance
team does not substitute a cheaper tool; it ships without systematic testing at
all. Measured cost per executed test is therefore a societal variable as well
as a budgetary one, and it is the reason NFR-D1 exists.
\emph{Who is affected:} small development firms, disproportionately in
economies like Bangladesh's, and the users of the software they ship.
\emph{How the design responds:} unit cost is measured rather than assumed
(DR-10), the expensive roles were moved to cheaper model tiers once the quality
cost of doing so had been measured, and a new target is configured as a
validated JSON profile through the dashboard rather than as a code change
(DR-09), so adoption does not require an engineer who can modify the system.

\paragraph{Accessibility, as a by-product.} The Android executor drives
applications through the platform accessibility tree and the web executor
through the accessibility-annotated DOM. The consequence is worth stating
plainly: a control this agent cannot locate is, in general, a control that a
screen-reader user cannot locate either. The agent's element-location failures
are therefore an accessibility signal, and its failure log doubles as a rough
accessibility audit of the application under test.
\emph{Who is affected:} users of assistive technology.
\emph{How the design responds:} controls are addressed by accessible role, name and content description rather than by pixel coordinates (\Cref{sec:standards}), and each agent-attributed failure records the specific control that could not be reached. The agent does also receive a screenshot, which sharpens the signal rather than weakening it: a control the agent can see in the image but cannot name from the accessibility tree is exactly a control that is visible to a sighted user and invisible to a screen reader, and that divergence is the accessibility gap itself. This is a genuine benefit but an incidental
one; it should not be presented as a designed accessibility audit, because the
agent does not check contrast, focus order, or any other WCAG criterion.

\paragraph{Language, and a failure that was silent.} The project's text
normalisation originally filtered input to the Latin range, which strips every
non-Latin character. Two Bengali strings therefore both normalised to the empty
string and were scored as nought per cent similar regardless of what they
actually said. Bengali vowel signs are Unicode combining marks and are not
covered by the usual word-character classes, so the bug survived casual
inspection. It was fixed by normalising on Unicode category rather than on a
Latin character class. The general point is larger than the bug: tooling built
on Latin-script assumptions does not fail loudly on a Bengali interface, it
fails silently and reports a plausible number.
\emph{Who is affected:} users and testers of software in Bengali and other
non-Latin scripts. \emph{How the design responds:} script-agnostic
normalisation wherever user-visible text is compared.

\paragraph{Legal considerations.} These are analysed in full in
\Cref{sec:standards}. In summary: dedicated disposable test accounts rather
than production user data, no ingestion of user records into the knowledge
graph, and destructive or communicative actions blocked at the configuration
layer rather than discouraged in prose.

\subsection*{Negative impacts, and what is done about them}

\begin{enumerate}[label=\textbf{N\arabic*.},leftmargin=3.0em]
  \item \textbf{Effect on quality-assurance employment.} A tool whose stated
    purpose is to displace manual effort has a labour effect, and it would be
    dishonest to present it otherwise. The evidence available suggests
    augmentation rather than replacement so far: the closest field deployment in
    the literature reports a saving of 1.5 person-days per regression round on a
    single business line~\cite{auitestagent}, which redirects effort rather than
    removing a role. The design choice that matters here is that the output is a
    natural-language test case a person can read, review and take ownership of,
    not an opaque verdict, so the human role shifts toward judgement rather than
    execution. At larger scale the effect on entry-level testing work is real
    and lies outside what this project can control.
  \item \textbf{Automation bias.} The most likely way this system causes harm is
    not by failing, but by being believed. An organisation that treats an
    automated verdict as authoritative will ship the defects the agent missed.
    \emph{Response:} a reported defect is explicitly scoped as a candidate for
    human confirmation and never as a confirmed fault (\Cref{sec:objectives});
    the generated report states its own degradations rather than hiding them
    (NFR-D2); and the attribution taxonomy makes it structurally impossible for
    the agent's own failure to be filed as an application defect (NFR-D3).
  \item \textbf{Acting on a production system with live users.} The evaluation
    target is in production. An exploratory agent that decides for itself what
    to do is, by construction, capable of doing something irreversible.
    \emph{Response:} the safety boundary is machine-enforced (DR-07) through
    blocked control texts and blocked URL patterns, out-of-scope requirements are
    withheld from the planner's citable list, and accounts used for testing are
    disposable.
  \item \textbf{Prompt injection.} The application's own on-screen text is fed
    to a model that also decides the next action, so content under an attacker's
    control could in principle direct the agent.
    \emph{Response:} every prompt that carries observed application content
    labels it explicitly as untrusted data that must never be read as
    instructions, in line with the OWASP guidance for language-model
    applications~\cite{owaspllm}.
  \item \textbf{False reports wasting developer time.} Each incorrectly reported
    defect costs investigation effort and erodes trust in the tool, which is the
    conflict recorded in \Cref{tab:stakeholders}.
    \emph{Response:} the evidence-integrity requirements above. The honest
    limitation is that the false-positive rate has not been measured, because
    doing so requires the seeded-defect experiment described in
    \Cref{sec:conclusion}.
\end{enumerate}

\section{Environmental Impact and Life Cycle Sustainability}
\label{sec:sustainability}

\subsection*{Development}

The development footprint is five developer workstations, one physical Android
handset, emulator instances on those workstations, and a local graph database.
No hardware was manufactured for the project; the handset is an existing
personal device. The decision that matters most environmentally was taken for
other reasons: model training and fine-tuning were placed out of scope
(\Cref{sec:objectives}), and the project uses hosted inference through an API
instead. Training is the dominant energy cost in most machine-learning work, so
excluding it means this project's footprint is an inference footprint only.

\subsection*{Operation}

The recurring environmental cost is hosted model inference, and it scales with
the number of tokens processed rather than with wall-clock time. Three
properties of the design reduce that volume. All three were adopted to control
money rather than energy, but the mechanism is the same.

\begin{itemize}[leftmargin=1.6em]
  \item Roughly nine tenths of the work in each execution happens on the device
    or in the browser rather than in a model, so most of a campaign consumes no
    inference at all.
  \item The executor and investigator roles run on smaller, cheaper model tiers,
    adopted once the quality cost of the downgrade had been measured.
  \item Context is budgeted rather than accumulated. Findings are atomic and
    deduplicated on write (DR-06), so a discovery made in round one is carried
    into round twenty as a single line with an observation count instead of
    being restated in full each time. Before that change, one test case's
    prose occupied the majority of the next generation prompt.
\end{itemize}

The literature supports token volume as the correct lever. AutoDroid halves its
prompt length by merging functionally equivalent interface nodes and reports
the query cost falling by the same proportion~\cite{autodroid}; AppAgentX
reduces tokens per task from 9.26k to 4.94k by replacing repeated reasoning
with stored shortcuts~\cite{appagentx}. Against these, the per-application
figures for a two-hour exploratory campaign reported by
DroidAgent~\cite{droidagent} indicate the scale of inference that an
unoptimised design consumes.

Device-side energy is smaller but not zero: a handset driven continuously for
the length of a campaign, plus a graph database left running between campaigns.
The honest limitation is that this project measured \emph{monetary} cost, not
energy. Token count is a reasonable proxy, but converting it to energy or
carbon requires the provider's hardware and grid mix, which are not disclosed.
No carbon figure is therefore claimed anywhere in this report.

\subsection*{Data life cycle, which is the real issue}

The knowledge graph accumulates screenshots, accessibility dumps and execution
traces of a production application, gathered through accounts that can sign in
to it. That is sensitive local test evidence with a retention question
attached, not neutral telemetry, and three concrete problems follow.

\begin{enumerate}[label=\textbf{E\arabic*.},leftmargin=3.0em]
  \item \textbf{Decay is not deletion.} A knowledge half-life (default 90 days)
    downweights stale knowledge in retrieval so the agent adapts when an
    application changes. It keeps the graph \emph{relevant}; it does nothing to
    bound how large it grows or how long the underlying evidence persists.
    \emph{Mitigation:} a retention policy that actually deletes, applied to
    screenshots and raw traces separately from the derived findings that give
    them value.
  \item \textbf{Graph snapshots are large, and permanent if mishandled.} A full export of the knowledge graph runs to tens of megabytes and carries the same evidence the live graph does. Version control is the wrong home for such a snapshot, because anything committed to a repository is effectively permanent in its history. \emph{Mitigation:} snapshots are held outside version control, under the same controls as the test credentials and with a stated retention period.
  \item \textbf{Screenshots are the bulk and the most sensitive part.}
    \emph{Mitigation:} store them outside the repository with an explicit
    lifetime, and keep the derived caption and structural signature, which is
    what retrieval actually uses, rather than the image, once that lifetime
    expires.
\end{enumerate}

\subsection*{Retirement}

The system is four processes, one database and no bespoke hardware, so
decommissioning it is deleting a database, revoking API keys and disabling the
test accounts. Nothing becomes electronic waste. The obligation that outlives
the software is the data described above, and it is the part that a retirement
plan has to name explicitly rather than assume will disappear with the code.

\section{Ethics}
\label{sec:ethics}

\subsection{Equity and Inclusivity}
\label{sec:equity}

Three groups could be excluded by a system of this kind, and each was
considered.

\emph{Users of software written in non-Latin scripts.} The normalisation defect
described in \Cref{sec:society} is the concrete instance: a component that
silently scored every pair of Bengali strings as completely dissimilar. It is
representative of a wider pattern in which tooling developed against English
interfaces degrades invisibly elsewhere, and the response was to normalise on
Unicode category so that any script is handled on the same footing.

\emph{Users of assistive technology.} Covered in \Cref{sec:society}: because
the agent is confined to what the accessibility layer exposes, an application
that is opaque to it is generally also opaque to a screen reader, and the
agent's failures make that visible rather than hiding it.

\emph{Organisations without the budget for quality assurance.} The measured
unit cost and the configuration-driven target system are what put this within
reach of a team that cannot fund a dedicated testing function.

The limitation to state honestly is that the agent is only as inclusive as the
interface it is given. An application built with a toolkit that exposes no
stable identifiers and no text through the accessibility tree, which is
precisely constraint C4, is one the agent tests poorly. That is a real
exclusion and it falls on exactly the applications that are already least
accessible.

\subsection{Accountability and Personal Responsibility}
\label{sec:accountability}

Responsibility for a system that issues correctness judgements has to be
located deliberately, because the failure mode is a machine-authored claim that
a human then acts on.

At the system level the position taken throughout this project is that
\textbf{the agent is not accountable for its verdicts; the people operating it
are}. Three mechanisms make that position operable rather than rhetorical. A
reported defect is scoped as a candidate requiring human confirmation
(\Cref{sec:objectives}). The three-way fault attribution taxonomy (NFR-D3) has
exactly one definition in the source tree, so an agent limitation cannot be
reported as an application defect through drift between components. And the
cross-process degradation register (NFR-D2) records every fallback taken, on
the principle that a system that hides its own degradation is worse than one
that fails loudly.

Two decisions during the work illustrate the same principle applied to the
team's own conduct. The team put its own codebase through a structured internal
review against an explicit severity rubric, rather than assuming that code it
had written was code it had checked, and scheduled the resulting work
accordingly. And the evaluation reports autonomy and coverage without
adjustment, including where those numbers are unflattering, because a testing
tool that presents its own results selectively has disqualified itself from the
job.

\guidance{TODO (team): add the per-member statement this indicator also asks
for, naming who owned which subsystem and who is answerable for it. That is the
same information as \Cref{sec:participation} and should be consistent with it.}

\subsection{Proper Use of Intellectual Property}
\label{sec:ip}

\paragraph{Prior work and literature.} Every empirical claim taken from the
research literature in \Cref{ch:research} is attributed to the paper it came
from, and the figures were read from those papers directly rather than from
secondary summaries. Where a work was surveyed rather than read in full, this
is stated at the point of use (\Cref{subsec:additional}) instead of being
blurred.

\paragraph{Third-party software.} The system is assembled from open-source
components and hosted services. The principal Python dependencies are FastAPI
and Uvicorn (HTTP services), Pydantic (schemas), the Neo4j driver, LangChain
and LangGraph (agent orchestration), Playwright (browser execution), mobilerun
(Android device execution), fastembed and NumPy (embeddings), markitdown, pypdf
and fpdf2 (document handling), structlog and psutil (operations), and an
OpenRouter client for hosted inference. The dashboard uses React built with
Vite. Neo4j itself is the graph database. Model inference is obtained
commercially through OpenRouter and is governed by that provider's terms of
service.

\guidance{TODO (team): before submission, list each of the above with its
licence (most are MIT, Apache-2.0 or BSD; Neo4j Community is GPLv3 and
Playwright is Apache-2.0) and confirm that the project's own intended licence
is compatible with all of them. \texttt{pip-licenses} will generate the table.}

\paragraph{The partner's material.} The technical requirements document
ETA-TR-2026-001 is not public. It is cited for traceability, and requirement
identifiers and titles are reproduced where they are needed to show what was
delivered, but no confidential content beyond that appears in this report.

\paragraph{Generative AI.} Generative AI tools were used during this project,
and the Candidates' Declaration requires that use to be disclosed rather than
concealed. Their role was assistance with implementation, refactoring,
documentation and drafting, under the team's direction and review; the
architecture, the requirements analysis, the experimental decisions and the
interpretation of results are the team's own. All generated code was read,
tested and, where wrong, corrected before it was kept.

\guidance{TODO (team): replace the sentence above with the team's own precise
account of how these tools were used, and check it against whatever disclosure
wording the department requires. Do not submit a generic statement.}

\subsection{Professionalism and Ethical Codes}
\label{sec:codes}

The ACM Code of Ethics and the IEEE Code of Ethics both place an obligation on
the engineer to avoid harm, to be honest about the limitations of their work,
and to accept responsibility for professional judgements. The IEB code makes
the same demand of the engineer's duty to the public. Three decisions in this
project were settled by those obligations rather than by convenience, and each
is documented at the point where it was taken.

\emph{Honesty about what the work establishes.} The review in
\Cref{subsec:measurement} concludes that this project cannot report precision or
recall for defect detection without seeding known faults, and says so, rather
than substituting a coverage number that would look like a result. The
corresponding ACM obligation is to be honest and trustworthy, including about
the limits of one's own competence and evidence.

\emph{Avoiding harm to a third party's users.} The evaluation target is a
production system with real users. Registration, one-time-password
verification, payment and account deletion were placed out of scope and blocked
in configuration rather than merely discouraged (DR-07), on the straightforward
reading that an engineer may not risk another party's users to make their own
evaluation more complete.

\emph{Applying to our own system the standard we ask of others.} The test account's secret is passed only to the component that has to type it, and never to the planner, which receives the account role alone (NFR-D6). The planner would produce marginally better-grounded tests with the full credential in its context, so the decision cost a little capability to uphold a principle that would have been invisible had it gone the other way. An engineer's obligations do not lapse because the system under discussion is their own.
